%DS01_ThanhNgân
\section{ĐỀ SỐ \thesection } 
\setcounter{bt}{0}
%%Câu 1.
\begin{bt} Thực hiện các phép tính
	\begin{listEX}[3]
		\item $130+(-12)+(-30)+12$;
		\item $25+75: 5-2.5^{2}$;
		\item $125+\left[-12-(7-3)^{3}\right]$.
	\end{listEX}
	
	\loigiai{
		\begin{listEX}
		\item $130 + (-12) + (-30) + 12 = 130 - 12 - 30 + 12 = 100$ ;
		\item $25 + 75 : 5 - 2 \cdot 5^{2} = 25 + 15 - 2 \cdot 25 = 25 + 15 - 50 = -10$ ;
		\item $125 + \left[-12 - (7 - 3)^{3}\right]
		= 125 + \left[-12 - 4^{3}\right] = 125 + \left[-12 - 64\right] = 125 - 76 = 49$.
	\end{listEX}
	}
\end{bt}


%%Câu 2.
\begin{bt} \textbf{ }
	\begin{listEX}
		\item Sắp xếp theo thứ tự tăng dần: $-29 ;\ 6 ;\ -30\,009 ;\  0 ;\  30$.
		\item Cho tập hợp $C=\{x \in \mathbb{Z} \mid-4 \leq x<3\}$. Hãy viết tập hợp $C$ bằng cách liệt kê các phần tử.
	\end{listEX}
	
	
	\loigiai{a) Thứ tự tăng dần là: $-30\,009 ;\ -29 ;\ 0 ;\ 6 ;\ 30$\qquad
		b) $C = \{-4;\  -3;\  -2;\  -1;\  0;\  1;\  2\}$
	}
\end{bt}


%%Câu 3.
\begin{bt} Tìm số nguyên $x$, biết
	\begin{listEX}[2]
		\item $18-3 x=3$;
		\item $18 \ \vdots \ x$ và $x>6$.
	\end{listEX}
	
	\loigiai{
	 \begin{listEX}[2]
		\item 18 - 3x = 3\\
		\hspace*{0.5 cm}$\begin{aligned}
			-3x &= 3 - 18\\
			-3x &= -15\\
			x &= 5
		\end{aligned}$ 
		\item $18\ \vdots \ x \Rightarrow x \in $ Ư$(18)=\{\pm1;\pm2;\pm3;\pm6;\pm9;\pm18\}$\\
			và $ x> 6$ nên $x \in\{9;18\}$ 
	\end{listEX}	 
	}
\end{bt}


%%Câu 4.
\begin{bt} \textbf{ }
	\begin{listEX}
		\item Lớp 6A có một số quyển tập cần gói thành những phần quà tặng cho học sinh có nhiều tiến bộ của lớp. Nếu mỗi phần quà có $8$ quyển, $12$ quyển hoặc $15$ quyển thì vừa đủ tập. Cho biết số quyển tập trong khoảng $300$ đến $400$ quyển. Tính số quyển tập của lớp 6A?
		\item Mỗi toa tàu chở được 20 tấn gạo. Hỏi nếu có 425 tấn gạo thì phải dùng bao nhiêu toa tàu?
	\end{listEX}
	
	\loigiai{
		\begin{listEX}
			\item Gọi số quyển tập là $a (a\in \mathbb{N})$. Ta có $a$ là bội chung của $8$, $12$, $15$ và $300 < a< 400$\\ 
			Phân tích mỗi số ra thừa số nguyên tố :
			$8=2^3$ \hspace{1cm} $12=2^2\cdot3$ \hspace{1cm} $15=3\cdot5$\\
			BCNN($8$, $12$, $15$)=$2^3\cdot3\cdot5=120$\\
			BC($8$, $12$, $15$)=B($120$)=\{$0$; $120$; $240$; $360$; $480$\}. Do $300 < a< 400$ nên $a=360$.\\
			Vậy số quyển tập của lớp 6A là $ 360$ quyển.
			\item Số toa tàu cần dùng là: $\frac{425}{20} = 21{,}25$.\\
			Vậy cần 22 toa để chở đủ 425 tấn gạo.
	\end{listEX}
		 
	}
\end{bt}


%%Câu 5.
\begin{bt} Một khu vườn hình chữ nhật có chiều dài $30$ m, chiều rộng ngắn hơn chiều dài $10$ m.
	\begin{listEX}[2]
		\item Tính chu vi khu vườn.
		\item Tính diện tích khu vườn.
	\end{listEX}
	
	\loigiai{a) Chiều rộng của khu vườn là: $30 - 10 = 20$ m.\\
		Chu vi khu vườn là: $2\cdot(30 + 20) = 2\cdot50 = 100$ m.\\
		b) Diện tích khu vườn là: $30\cdot 20 = 600$ m$^2$.
	}
\end{bt}


%%Câu 6.
\begin{bt} \textbf{ }
	\begin{listEX}
		\item Một chiếc diều bay lên đến độ cao $15$ m (so với mặt đất), sau đó giảm $5$ m rồi gặp gió lại bay lên $7$ m. Hỏi cuối cùng chiếc diều ở độ cao bao nhiêu (so với mặt đất)?
		\item Chứng minh: $1+5+5^{2}+5^{3}+5^{4}+5^{5}+5^{6}+5^{7}+5^{8}$  chia hết cho $31$.
	\end{listEX}
	
	
	\loigiai{
		\begin{listEX}
		\item Độ cao cuối cùng của chiếc diều là: $15 - 5 + 7 = 17$ m.
		\item  $1 + 5 + 5^{2} + 5^{3} + 5^{4} + 5^{5} + 5^{6} + 5^{7} + 5^{8}\\
		=(1 + 5 + 5^{2}) + 5^{3}\cdot(1 + 5 + 5^{2}) + 5^{6}\cdot(1 + 5 + 5^{2})
		= 31 + 5^{3}\cdot31 + 5^{6}\cdot31 \\
		= 31\cdot(1+5^{3}+5^{6}) $ do $31\vdots 31$ nên tích $31\cdot(1+5^{3}+5^{6})$ chia hết cho $31$\\
		Vậy $1 + 5 + 5^{2} + 5^{3} + 5^{4} + 5^{5} + 5^{6} + 5^{7} + 5^{8}$ chia hết cho 31
	 \end{listEX}}
\end{bt}

